%%%%%%%%%%%%%%%%%%%%%%%%%%%%%%%%%%%%%%%%%%%%%%%%%%%%%%%%%%%%%%%%%%%%%%
% Lớp 10 | Hàm số bậc hai và đồ thị | Hàm số và đồ thị
% Dạng 0D3?1-1: Xác định một hàm số
% Số câu: 30  (sinh tự động bởi tach_cau.py; từ nay thầy sửa trực tiếp file này)
%%%%%%%%%%%%%%%%%%%%%%%%%%%%%%%%%%%%%%%%%%%%%%%%%%%%%%%%%%%%%%%%%%%%%%

%%%%%%%%%%%%%%%% TRẮC NGHIỆM 4 PHƯƠNG ÁN %%%%%%%%%%%%%%%%

% nguồn: MaHoaNH_TracNghiem_Lop10
\begin{ex}%[0D3N1-1]%[ma:ecc50d2c]
	 Hàm số nào dưới đây là hàm số bậc nhất?
	\choice 
		 {$y=2x+\dfrac{1}{x}$}
		 {$y=2$}
		 {$y=\dfrac{x+1}{x-2}$}
		{\True $y=2x+\sqrt{2}$} 
	 \loigiai{Theo định nghĩa về hàm số bậc nhất, $y=2x+\sqrt{2}$ là hàm số bậc nhất.}
\end{ex}

% nguồn: MaHoaNH_TracNghiem_Lop10
\begin{ex}%[0D3N1-1]%[ma:363d450d]
	Biết rằng đồ  thị hàm số $y=ax+b$ đi qua điểm $A(-3;1)$ và có hệ số góc bằng $-2$. Tính tích $P=ab$.
	 \choice
		 {$P=-10$}
		 {\True $P=10$}
		{$P=-7$} 
		 {$P=7$}
	\loigiai{ 
		Do đồ thị hàm số $y=ax+b$ đi qua điểm $A(-3;1)$ và có hệ số góc bằng $-2$ ta  có hệ phương trình
	$$\heva{&a=-2\\ &1=-3a+b} \Leftrightarrow \heva{&a=-2\\ &b=-5}  \Rightarrow P=(-2) \cdot (-5)=10.$$}
\end{ex}

% nguồn: MaHoaNH_TracNghiem_Lop10
\begin{ex}%[0D3N1-1]%[ma:5714e431]
	Trong các hàm  số sau, hàm số nào là hàm số bậc nhất?
	\choice 
		{$y=(-4m+2)x+7$} 
		{$y=2x+3x^2$} 
		 {\True $y=10x+5$}
		{$y=9$} 
	\loigiai{ 
		Hàm số bậc nhất là hàm số  có dạng $y=ax+b$ với $a\ne 0$.\\
	Vậy hàm số $y=10x+5$ là hàm  số bậc nhất.}
\end{ex}

% nguồn: MaHoaNH_TracNghiem_Lop10
\begin{ex}%[0D3N1-1]%[ma:ea269b0d]
	Đồ  thị của hàm số $y=f(x)=\heva{&2x+1 &\text{khi} && x\le 2\\&-3 &\text{khi} && x>2}$ đi qua điểm nào sau đây?
	 \choice
		 {$(0;-3)$}
		 {$(3;7)$}
		 {$(2;-3)$}
		{\True  $(0;1)$}
	\loigiai{ 
		Thử lần lượt từng phương án A, B, C, D  với chú ý về điều kiện ta được:\\
		$f(0)=2\cdot0+1=1 \ne-3$, đồ thị không đi  qua điểm $(0;-3)$.\\
		$f(3)=-3\ne 7$, đồ  thị không đi qua điểm $(3;7)$.\\
		$f(2)=2\cdot2+1=5\ne-3$, đồ  thị không đi qua điểm $(2;-3)$.\\
	$f(0)=2\cdot 0+1=1$,  đồ thị đi qua điểm $(0;1)$.}
\end{ex}

% nguồn: MaHoaNH_TracNghiem_Lop10
\begin{ex}%[0D3H1-1]%[ma:864b0e56]
	Cho hàm số bậc nhất $y=ax+b$. Tìm $a$  và $b$, biết rằng đồ thị hàm số đi qua điểm $M(-1;1)$ và cắt trục hoành tại điểm có hoành độ là $5$.\\
	\choice 
		{$a=\dfrac{1}{6}$; $b=\dfrac{5}{6}$} 
		{$a=-\dfrac{1}{6}$; $b=-\dfrac{5}{6}$} 
		{$a=\dfrac{1}{6}$; $b=-\dfrac{5}{6}$} 
		{\True  $a=-\dfrac{1}{6}$; $b=\dfrac{5}{6}$}
	\loigiai{ 
	Theo  đề ta có $\heva{&1=-a+b\\ &0=5a+b} \Leftrightarrow \heva{&a=-\dfrac{1}{6}\\ &b=\dfrac{5}{6}.}$}
\end{ex}

% nguồn: MaHoaNH_TracNghiem_Lop10
\begin{ex}%[0D3H1-1]%[ma:8ca03409]
	Cho đường thẳng có phương trình $y=ax+b$ đường thẳng đi qua  hai điểm $M(1;3)$, $N(2;-4)$. Giá trị của $a$ và $b$ là
	\choice 
		{$a=-7$;  $b=-10$}
		 {$a=7$; $b=10$}
		{$a=7$; $b=-10$} 
		{\True  $a=-7$; $b=10$}
	\loigiai{ 
		Đường thẳng đã cho đi qua $M(1;3)$, $N(2;-4)$ nên ta có hệ phương  trình\\
	$$\heva{&a+b=3\\ &2a+b=-4}  \Leftrightarrow \heva{&a=-7\\ &b=10.}$$}
\end{ex}

% nguồn: MaHoaNH_TracNghiem_Lop10
\begin{ex}%[0D3H1-1]%[ma:4ea9f4b9]
	Tìm $a$ và $b$ biết rằng đường thẳng $y=ax+b$ đi qua $M(1;-1)$ và song song với  đường $y=2x+3$.
	\choice 
		{$\heva{&a=-1\\  &b=2}$}
		{\True $\heva{&a=2\\  &b=-3}$}
		{$\heva{&a=2\\  &b=4}$}
		{$\heva{&a=2\\  &b=3}$}
	\loigiai{ 
		Đường thẳng $y=ax+b$ đi qua $M(1;-1)$ và song song với đường $y=2x+3$ nên ta có\\ 
	$$\heva{&a+b=-1\\  &a=2\\ &b\ne 3} \Leftrightarrow \heva{&a=2\\ &b=-3.}$$}
\end{ex}

% nguồn: MaHoaNH_TracNghiem_Lop10
\begin{ex}%[0D3H1-1]%[ma:dd822eaf]
	Tập xác định của hàm  số $y=\sqrt{x^{2}-1}$ là\\
	 \choice
		 {$\mathscr{D}=(-\infty;0]$}
		{$\mathscr{D}=(-1;1)$} 
		{$\mathscr{D}=[0;+\infty)$} 
		 {\True $\mathscr{D}=(-\infty;-1]\cup[1;+\infty)$}
	\loigiai{ 
		Hàm số $y=\sqrt{x^{2}-1}$ xác định khi $x^{2}-1\geq  0\Leftrightarrow\hoac{&x\geq 1\\&x\leq-1.}$\\
		Vậy hàm số $y=\sqrt{x^{2}-1}$ có tập xác định  là $D=(-\infty;-1]\cup[1;+\infty)$.\\
	 }
\end{ex}

% nguồn: MaHoaNH_TracNghiem_Lop10
\begin{ex}%[0D3H1-1]%[ma:1efa1431]
	 Biết rằng $(P)\colon y=ax^2-4x+c$ có hoành độ đỉnh bằng $-3$ và đi qua điểm $M(-2;1)$. Tính tổng $S=a+c$.
	 \choice
		{$S=5$} 
		 {\True $S=-5$}
		{$S=4$} 
		 {$S=1$}
	\loigiai{ 
		Ta có: 
		$\dfrac{-b}{2a}=-3\Leftrightarrow\dfrac{2}{a}=-3\Leftrightarrow a=-\dfrac{2}{3}$ và $4a+8+c=1 \Leftrightarrow c=-\dfrac{13}{3}$.\\ 
		 Vậy $S=a+c=-5$.
	 }
\end{ex}

% nguồn: MaHoaNH_TracNghiem_Lop10
\begin{ex}%[0D3H1-1]%[ma:ec8d50b7]
	Cho hai hàm số $y=(m+1)^2x-2$ và $y=(3m+7)x+m$. Tìm tất cả các giá trị thực của tham số $m$ để đồ thị hai hàm số đã  cho cắt nhau.
	 \choice
		{$m\ne-2$} 
		 {$m\ne-3$}
		{\True $m\ne-2$,  $m\ne 3$}
		{$m=-2$,  $m=3$}
	\loigiai{ 
		Hoành độ  giao điểm của đồ thị hai hàm số đã cho là nghiệm của phương trình\\
		$$(3m+7)x+m=(m+1)^2x-2 \Leftrightarrow (m^2-m-6)x=m+2.  \qquad(*)$$
		Đồ thị hai hàm  số đã cho cắt nhau khi và chỉ khi (*) có $1$ nghiệm
	$$\Leftrightarrow  m^2-m-6\ne 0 \Leftrightarrow \heva{&m\ne 3\\ &m\ne-2.}$$}
\end{ex}

% nguồn: MaHoaNH_TracNghiem_Lop10
\begin{ex}%[0D3H1-1]%[ma:e81bf132]
	Cho hàm số bậc nhất $y=ax+b$ có đồ thị là đường thẳng  đi qua điểm $A(0;-1)$ và song song với đường thẳng $y=3-2x$. Tính tổng $S=2a+b$.
	\choice 
		 {\True $-5$}
		 {$4$}
		 {$6$}
		 {$-3$}
	 \loigiai{
		Vì $\mathrm{d} \colon y=ax+b$ song song với đường  thẳng $y=3-2x$ nên $a=-2$ và $b\ne 3$.\\
		Suy  ra $d \colon y=-2x+b \qquad(b\ne 3)$.\\
		Vì $\mathrm{d} \colon y=-2x+b$ đi qua $A(0;-1)$ nên $-1=-2 \cdot 0+b  \Rightarrow b=-1$.\\
	Vậy  $S=2a+b=2 \cdot (-2)+(-1)=-5$.}
\end{ex}

% nguồn: MaHoaNH_TracNghiem_Lop10
\begin{ex}%[0D3H1-1]%[ma:a8474314]
	Cho ba đường thẳng $d_1\colon y=x+7$, $d_2 \colon y=ax-4$, $d_3 \colon y=2x+8$. Biết rằng $d_1$, $d_2$, $d_3$ đồng qui tại điểm $M(b;c)$.  Tính tổng $S=a+b+c$.
	\choice 
		 {$-1$}
		 {\True $-5$}
		{$-10$} 
		{$6$} 
	\loigiai{ 
		Phương trình hoành độ giao điểm của $d_1$,  $d_3$ là $x+7=2x+8 \Leftrightarrow x=-1$.\\
		Thay $x=-1$ vào $d_1 \colon y=x+7$, ta được  $y=6$.\\
		Vậy tọa độ  giao điểm của $d_1$, $d_3$ là $M(-1;6)$.\\
		Vì $d_1$, $d_2$, $d_3$ đồng  qui tại điểm $M(-1;6)$ nên $d_2$ đi qua $M(-1;6)$.\\
		 Khi đó $6=a(-1)-4\Rightarrow a=-10$.\\
	 Vậy $S=a+b+c=-10+(-1)+6=-5$.}
\end{ex}

% nguồn: MaHoaNH_TracNghiem_Lop10
\begin{ex}%[0D3H1-1]%[ma:80c32056]
	Giá trị  của tham số $m$ để ba đường thẳng $y=x+1$, $y=-x+m$, $y=3x$ đồng quy tại một điểm là
	 \choice
		{$m=1$} 
		 {\True $m=2$}
		 {$m=\dfrac{-1}{2}$}
		 {$m=\dfrac{1}{2}$}
	 \loigiai{
		Giả sử $d_1 \colon y=x+1$, $d_2  \colon y=-x+m$, $d_3 \colon y=3x$.\\
		Gọi $A$ là giao điểm của $d_1$  và $d_3$.\\
		Phương trình hoành độ  giao điểm của $d_1$ và $d_3$ là
		$$x+1=3x \Leftrightarrow x=\dfrac{1}{2} \Rightarrow  y=\dfrac{3}{2}.$$\\
		Suy ra  $A\left(\dfrac{1}{2};\dfrac{3}{2}\right)$.\\
		Ba đường thẳng  $d_1$, $d_2$, $d_3$ đồng quy
		$$\Leftrightarrow A\in{d_2}  \Leftrightarrow\dfrac{3}{2}=-\dfrac{1}{2}+m \Leftrightarrow m=2.$$
	Vậy ba đường thẳng đã  cho đồng quy khi $m=2$.}
\end{ex}

% nguồn: MaHoaNH_TracNghiem_Lop10
\begin{ex}%[0D3V1-1]%[ma:56a4ab5b]
	Cho phương trình đường thẳng $y=ax+b$ có đồ thị đi qua điểm $E(2;-1)$ và song song với đường thẳng $ON$ với $O$ là gốc tọa độ  và $N(1;3)$. Tính giá trị biểu thức $S=a^2+b^2$.
	\choice 
		 {$S=-4$}
		{$S=-40$} 
		{$S=-58$} 
		{\True $S=58$} 
	\loigiai{ 
		Phương trình đường thẳng $ON$ với $O$ là gốc tọa độ và $N(1;3)$ có dạng  $y=3x$.\\
		 Đường thẳng $y=ax+b$ đi qua điểm $E(2;-1) \Leftrightarrow 2a+b=-1$. \qquad(1)\\
		Đường thẳng $y=ax+b$ song song với đường thẳng $ON \Rightarrow \heva{&a=3\\ &b\ne 0.}$  \qquad(2)\\
		Từ $(1)$ và $(2)$ suy  ra $\heva{&a=3\\ &b=-7.}$\\
	Vậy  $S=a^2+b^2=58$.}
\end{ex}

% nguồn: MaHoaNH_TracNghiem_Lop10
\begin{ex}%[0D3V1-1]%[ma:f079b63d]
	Cho hai đường  thẳng $d_1 \colon y=mx-4$ và $d_2 \colon y=-mx-4$. Gọi $S$ là tập hợp các giá trị nguyên dương của $m$ để tam giác tạo thành bởi $d_1$, $d_2$ và trục hoành có diện tích lớn hơn $8$. Số phần tử của tập $S$ là
	 \choice
		{\True $1$} 
		 {$3$}
		 {$2$}
		{$4$} 
	\loigiai{ 
		Ta có $d_1 \colon  y=mx-4$, $d_2 \colon y=-mx-4$ cắt nhau và cùng cắt trục $Ox$ khi $m\ne 0$.\\
		Gọi $A\left(\dfrac{4}{m};  0\right)$, $B\left(-\dfrac{4}{m}; 0\right)$ lần lượt là giao điểm của $d_1$, $d_2$ với trục hoành.\\
		Phương trình hoành độ giao  điểm của $d_1$, $d_2$ là $mx-4=-mx-4 \Leftrightarrow x=0$.\\
		Gọi $C$ là giao  điểm của $d_1;\,d_2$ thì $C(0;-4)$.\\
		Ta có $S_{\Delta ABC}=\dfrac{1}{2} \mathrm{d}(C,Ox) \cdot AB=\dfrac{1}{2}|y_C| \cdot |x_A-x_B|=\dfrac{1}{2} \cdot 4 \cdot \dfrac{8}{|m|}=\dfrac{16}{|m|}$.\\ 
		Có $S_{\Delta ABC}> 8 \Leftrightarrow\dfrac{16}{|m|}> 8\Leftrightarrow|m| < 2,\,m\in{N^*}  \Leftrightarrow m=1$.\\
	Vậy $S=\{ 1\}$.} 
\end{ex}

% nguồn: MaHoaNH_TracNghiem_Lop10
\begin{ex}%[0D3Y1-1]%[ma:23c2fab5]
	Từ một hộp chứa $5$ quả cầu trắng, $6$ quả cầu đỏ, $7$ quả cầu vàng, các quả cầu có kích thước và khối lượng giống nhau. Xét phép thử: "Lấy ngẫu nhiên đồng thời $3$ quả cầu". Biến cố đối của biến cố $A\colon$ "$3$ quả cầu lấy ra có ít nhất một quả màu đỏ”  là
	\choice 
		{$\bar{A}\colon$ "3 quả  cầu lấy ra cùng màu trắng"}
		{\True $\bar{A}\colon$ "3 quả cầu lấy ra  không có màu đỏ"}
		{$\bar{A}\colon$ "3 quả  cầu lấy ra cùng màu vàng"}
		{$\bar{A}\colon$ "3 quả cầu lấy ra có  3 màu khác nhau"}
	\loigiai{ 
		Biến cố đối của biến cố $A$ là $\bar{A}\colon$ "3 quả cầu lấy ra không có  màu đỏ".
	} 
\end{ex}

% nguồn: MaHoaNH_TracNghiem_Lop10
\begin{ex}%[0D3B1-1]%[ma:4cbfdbf5]
	Gieo một con xúc xắc hai lần.  Tập $\big\{(1;3),\;(2;4),\;(3;5),\;(4;6)\big\}$ là biến cố nào dưới đây?
	\choice 
		{$N$:  ``Tổng số chấm hai lần gieo là chẵn''}
		{$Q$: ``Số chấm  hai lần gieo hơn kém $2$''}
		{\True $M$: ``Lần thứ hai hơn lần  thứ nhất hai chấm''}
		{$P$: ``Tích số  chấm hai lần gieo là chẵn''}
	 \loigiai{
		$\big\{(1;3),\;(2;4),\;(3;5),\;(4;6)\big\}$  là biến cố $M$: ``Lần thứ hai hơn lần thứ nhất hai chấm''
	} 
\end{ex}

% nguồn: MaHoaNH_TracNghiem_Lop10
\begin{ex}%[0D3Y1-1]%[ma:f33a2c49]
	Gieo $2$ đồng tiền là một phép thử  ngẫu nhiên có không gian mẫu là
	\choice 
		 {\True $\{NN,NS,SN,SS\}$}
		 {$\{NS,SN\}$}
		 {$\{N,S\}$}
		 {$\{NNN,SSS,NNS,SSN,NSS,SNN\}$}
	\loigiai{ 
		Ta có không gian mẫu là  $\Omega=\{NN,NS,SN,SS\}$.
	} 
\end{ex}

% nguồn: MaHoaNH_TracNghiem_Lop10
\begin{ex}%[0D3Y1-1]%[ma:f2b2f671]
	Gieo một con súc sắc, số phần tử  của không gian mẫu là
	 \choice
		 {$24$}
		 {$12$}
		{$8$} 
		{\True $6$} 
	\loigiai{ 
		Ta có $n(\Omega)=6$. 
	 }
\end{ex}

% nguồn: MaHoaNH_TracNghiem_Lop10
\begin{ex}%[0D3Y1-1]%[ma:caa9b60e]
	Gieo một đồng xu cân đối và đồng chất 2 lần. Số phần tử của không gian mẫu là 
	\choice 
		{$  8 $}
		{$  6 $}
		{\True  $ 4 $}
		{$ 2 $ } 
	 \loigiai{
		Số  phần tử không gian mẫu $ n\left(\Omega\right) =4$.
	 }
\end{ex}

% nguồn: MaHoaNH_TracNghiem_Lop10
\begin{ex}%[0D3Y1-1]%[ma:dfa428ca]
	Một hộp có 5 quả cầu được đánh số từ $ 1 $ đến $ 5 $. Lấy ngẫu nhiên từ hộp đó liên tiếp $ 2 $ lần,  mỗi lần một quả và xếp thành hàng ngang. Số phần tử của không gian mẫu là
	 \choice
		{\True $ 20 $} 
		{$ 10  $}
		{$ 120 $} 
		{$  2 $}
	\loigiai{ 
		Số phần  tử của không gian mẫu là $ n\left(\Omega\right) =\mathrm{A}^2_5=20 $.
	 }
\end{ex}

% nguồn: MaHoaNH_TracNghiem_Lop10
\begin{ex}%[0D3Y1-1]%[ma:5813165d]
	Biểu thức nào  sau đây \textbf{không} là hàm số theo biến $x$?
	 \choice
		{$y=5x^2-3x+4$} 
		{$y=\sqrt{x^2-1}$} 
		{$y=\left|2x+3\right|$} 
		{\True $y^4=x^3$} 
	 \loigiai{
		Hàm số $y^4=x^3$ không phải là hàm số theo biến  $x$ vì ứng với giá trị của $x=1$ có hai giá trị của $y$ là $y=1$ và $y=-1$.
	} 
\end{ex}

% nguồn: MaHoaNH_TracNghiem_Lop10
\begin{ex}%[0D3Y1-1]%[ma:b80bbe70]
	Trong một chiếc hộp đựng $6$ viên bi đỏ, $8$ viên bi xanh, $10$  viên bi trắng. Lấy ngẫu nhiên $4$ viên bi.
	Số phần  tử của không gian mẫu là?
	 \choice
		 {\True $n(\Omega)=\mathrm{C}_{24}^4$}
		 {$n(\Omega)=\mathrm{A}_{24}^4$}
		 {$n(\Omega)=\mathrm{C}_6^4+\mathrm{C}_8^4+\mathrm{C}_{10}^4$}
		{$n(\Omega)=6\cdot 8\cdot  10$}
	 \loigiai{
		Ta có  $n(\Omega)=\mathrm{C}_{24}^4$.
	 }
\end{ex}

% nguồn: MaHoaNH_TracNghiem_Lop10
\begin{ex}%[0D3Y1-1]%[ma:f7e13597]
	Có $100$ tấm thẻ được đánh số từ  $1$ đến $100$. Lấy ngẫu nhiên $5$ thẻ. Số phần tử của không gian mẫu
	 là
	\choice 
		 {\True $n(\Omega)=\mathrm{C}_{100}^5$}
		{$n(\Omega)=\mathrm{A}_{100}^5$} 
		 {$n(\Omega)=\mathrm{C}_{100}^1$}
		{$n(\Omega)=\mathrm{A}_{100}^1$} 
	\loigiai{ 
		Ta có $n(\Omega)=\mathrm{C}_{100}^5$. 
	} 
\end{ex}

% nguồn: MaHoaNH_TracNghiem_Lop10
\begin{ex}%[0D3Y1-1]%[ma:e39e3dd9]
	Tung một đồng xu cân đối và đồng chất hai lần. Số  phần tử của biến cố cả hai lần tung đều xuất hiện mặt ngửa là
	\choice 
		 {$2$}
		 {$3$}
		{\True $1$} 
		 {$4$}
	 \loigiai{
		Ta có  $n\left(\Omega\right)=2\cdot2=4$.\\
		Gọi $A$  là biến cố \lq\lq Cả hai lần tung đều xuất hiện mặt ngửa \rq\rq.\\
		 Suy ra $n\left(A\right)=1$.
	} 
\end{ex}

% nguồn: MaHoaNH_TracNghiem_Lop10
\begin{ex}%[0D3Y1-1]%[ma:d248ca34]
	Hàm số  mô tả sự phụ thuộc của $y$ (số tiền phải trả, tính bằng nghìn đồng) vào $x$ (lượng điện tiêu thu tính bằng kWh) trên từng khoảng giá trị $x$, được cho bởi công thức sau:
	$$ 
	y= \begin{cases}1{,}678 x & \text{Nếu}\; 0 \leq x \leq 50 \\ 1{,}734 x-2{,}8 & \text{Nếu}\; 0<x \leq 100 \\ 2{,}014 x-30{,}8 &\text{Nếu}\; 100<x \leq 200 \\ 2{,}536 x-135{,}2 &\text{Nếu}\; 200<x \leq 300 \\ 2{,}834 x-224{,}6 &\text{Nếu}\; 300<x \leq 400 \\ 2{,}927 x-261{,}8 & \text{Nếu}\; x>400
\end{cases}
$$ 
Ông An tháng thứ nhất sài hết $90$kWh điện, tháng thứ hai sài hết $120$kWh điện. Theo công thức trên số tiền ông An phải trà sau hai tháng gần nhất với kết quả nào  sau đây?
\choice 
	{\True $364$ nghìn  đồng}
	{$361$ nghìn đồng} 
	 {$392$ nghìn đồng}
	 {$352$ nghìn đồng}
 \loigiai{
	Tháng thứ nhất sài hết $90$kWh nên số tiền phải trà là: $1{,}734\cdot90-2{,}8=153{,}26$ nghìn  đồng.\\
	Tháng hai sài hết $120$kWh nên số tiền phải trà là: $2{,}014\cdot120-30{,}8=210{,}88$  nghìn đồng.\\
	Vậy  ông An phải trả sau hai tháng là: $153{,}26+210{,}88=364{,}14$ nghìn đồng.
 }
\end{ex}


%%%%%%%%%%%%%%%% ĐÚNG/SAI %%%%%%%%%%%%%%%%

% nguồn: MaHoaNH_DungSai_Lop10
\begin{ex}%[0D3N1-1]%[ma:72a96e1d]
	Cho  hàm số $y=f(x)$ với bảng giá trị tương ứng cho trong bảng sau
	\begin{center}
		\begin{tabular}{|c|c|c|c|c|c|c|c|c|c|}
			\hline 
			$x$ & $-5$ & $-3$ & $-1$ & $0$ &  $1$ & $2$ & $5$ & $8$ & $9$ \\
			 \hline
			$y$ & $-10$ & $-6$ & $-2$ & $0$ &  $2$ & $4$ & $8$ & $12$ & $15$ \\
			 \hline
		\end{tabular}
	\end{center}
	Mỗi  kết quả dưới đây đúng hay sai?
	 \choiceTF
		{\True Tập xác định  là $\left\{-5;-3;-1;0;1;2;5;8;9\right\}$}
		{$f(1)+f(-3)=0$} 
		{$y=2x$} 
		{\True Tập  giá trị là $\left\{-10;-6;-2;0;2;4;8;12;15\right\}$}
	 \loigiai{
		\begin{itemchoice}
			\itemch Từ giá trị của $x$ trong bảng ta có tập xác định là $\left\{-5;-3;-1;0;1;2;5;8;9\right\}$. 
			\itemch $f(-3)+f(1)=-6+2=-4$. 
			 \itemch Do $f(5)=8$ nên $y\ne 2x$.
			\itemch  Từ giá trị của của $y$ ta có tập giá trị là $\left\{-10;-6;-2;0;2;4;8;12;15\right\}$.
		\end{itemchoice}
	 }
\end{ex}

% nguồn: MaHoaNH_DungSai_Lop10
\begin{ex}%[0D3N1-1]%[ma:cc0b6702]
	Cho hàm số $y=f(x)$ với bảng giá trị tương  ứng cho trong bảng sau
	\begin{center}
		\begin{tabular}{|c|c|c|c|c|c|c|c|c|c|}
			 \hline
			$x$ & $0$ & $1$ & $2$ &  $3$ & $4$ & $5$ & $6$ & $7$ & $8$ \\
			\hline 
			$y$  & $1$ & $2$ & $3$ & $4$ & $5$ & $6$ & $7$ & $8$ & $9$ \\
			 \hline
		\end{tabular}
	\end{center}
	Mỗi kết quả dưới đây đúng  hay sai?
	 \choiceTF
		{Tập xác  định là $\mathscr{D}=\mathbb{N}$}
		 {\True $f(3)=4$}
		{\True  $y=x+1$}
		{Tập giá trị là $\left\{1;2;3;4;5;6;7;8;9;10\right\}$} 
	 \loigiai{
		\begin{itemchoice}
			\itemch Từ giá trị của $x$ trong bảng ta có tập xác  định là $\left\{0;1;2;3;4;5;6;7;8\right\}$.
			\itemch  $f(3)=4$.
			\itemch Từ bảng giá trị ta có  $y=x+1$.
			\itemch Từ giá trị của  của $y$ ta có tập giá trị là $\left\{1;2;3;4;5;6;7;8;9\right\}$.
		\end{itemchoice}
	 }
\end{ex}


%%%%%%%%%%%%%%%% TRẢ LỜI NGẮN %%%%%%%%%%%%%%%%

% nguồn: MaHoaNH_TraLoiNgan_Lop10
\begin{bt}%[0D3H1-1]%[ma:b2187c85]
	Cho hàm số $y=f(x)=-3x^2+m^2x+m+1$ (với $m$ là tham số). Có bao nhiêu giá trị của $m$ để đồ thị của hàm  số $y=f(x)$ đi qua điểm $A(1;0)$.
	\shortans{2} 
	 \loigiai{
		Đồ  thị của hàm số $y=f(x)$ đi qua điểm $A(1;0)$ khi và chỉ khi
		$$0=-3+m^2+m+1 \Leftrightarrow m^2+m-2=0  \Leftrightarrow \hoac{&m=1\\ &m=-2.}$$\\
	Vậy có $2$ giá trị của tham số $m$ thỏa mãn yêu  cầu bài toán.}
\end{bt}

% nguồn: MaHoaNH_TraLoiNgan_Lop10
\begin{bt}%[0D3H1-1]%[ma:72a3c616]
	Tổng  các số nguyên $m$ thuộc khoảng $(-3;3)$ để hàm số $y=\dfrac{2x+\sqrt 2}{x^2-2x+1-m}$, ($m$ là tham số) có tập xác định là $\mathbb{R}$.
	\shortans{-3} 
	\loigiai{ 
		Để hàm số $y=\dfrac{2x+\sqrt 2}{x^2-2x+1-m}$ có tập xác định  là $\mathbb{R}$ khi và chỉ khi phương trình $x^2-2x+1-m=0$ vô nghiệm hay $\Delta'=1-(1-m) < 0 \Leftrightarrow m < 0$.\\
		Với  $m$ thuộc khoảng $(-3;3)$ ta có $m=\{-2;-1\}$.\\
	Vậy tổng các số $m$ là  $-3$.}
\end{bt}
